\documentclass[10pt,a4paper]{amsart}
\usepackage[margin=19mm]{geometry}
\usepackage{amsmath,amssymb,mathtools}
\usepackage[scr=rsfs,cal=euler]{mathalfa}
\usepackage{xfrac}
\usepackage[unicode,hidelinks,bookmarks=false]{hyperref}
\newcommand{\ie}{i.e.\ }
\newcommand{\cf}{cf.\ }
\newcommand{\resp}{respectively\ }
\newcommand{\Subsection}{\subsection}
\providecommand{\nobreakdash}{\nobreak\hbox{-}}
\setlength{\emergencystretch}{2em}
\multlinegap0pt
\usepackage{bm}
\usepackage{bbm}
\usepackage{dsfont}
\usepackage{xspace}
\usepackage{cancel}
\usepackage{mathtools}
\usepackage{amsthm}
\makeatletter
\@ifundefined{lemma}{\newtheorem{lemma}{Lemma}}{}
\makeatother
\title[On the Gauss circle problem over smooth numbers]{On the Gauss circle problem over smooth numbers}
\author{Peng Gao}
\date{Source: arXiv:2604.23918v1 (CC BY 4.0)}
\begin{document}
\maketitle

\noindent\emph{Source and licence.} On the Gauss circle problem over smooth numbers. Version arXiv:2604.23918v1 (CC BY 4.0), distributed under the Creative Commons Attribution 4.0 International licence, as recorded in the arXiv licence metadata for this exact version and shown on the abstract page https://arxiv.org/abs/2604.23918v1. Full rights notice: licenses/arxiv-2604.23918-CC-BY-4.0.txt.\par\smallskip

\noindent\textit{Editorial scope.} Counted unit: the Lemma whose source label is \texttt{lemphibound} in the article (source0.tex lines 687--694, source label \texttt{lemphibound}), delivered together with the article's own complete original proof of it (source0.tex lines 695--795, one \texttt{proof} environment of 101 lines). No mathematics is written, repaired or re-derived: statement and proof are the article's own text line for line. Required context retained uncounted: sumpsigmalogp (lines 303--320); sumpsigmalogpchi (lines 303--320); alphaupper (lines 432--443); alphausmall (lines 432--443); alpharel (lines 447--452); sumpvy (lines 501--506); vest (lines 563--568); sumpvy4 (lines 577--582). Every definition, display and label of those lines is retained unchanged. Omitted from this package and not redistributed: 1--302, 321--431, 444--446, 453--500, 507--562, 569--576, 583--686, 796--1033. Nothing inside a retained excerpt is omitted or altered: every retained body is delivered exactly as the article's lines are written, so no omission receipt of any kind accompanies this package. Citations: the retained excerpts carry no citation command at all, so no bibliography entry of the article is transcribed and no reference is redistributed. Reference adaptation: none was needed and none was made; every reference inside the delivered unit resolves inside it, so no unresolved reference or citation is produced. Printed slips kept literal: no printed word, symbol, hypothesis or number of the counted statement or of its proof was altered. Layout and class: the article's own class is \texttt{amsart} with packages including amsfonts, amsmath, amscd, amssymb, amsthm, amsrefs, latexsym, mathrsfs, bbm, enumerate, graphicx, color; the wrapper is the lane's pinned \texttt{amsart} editorial wrapper at 10pt on A4, which restates the theorem-like environments the retained text needs and adds the article's own macro definitions that the retained text needs (none), with extra wrapper packages (bm, bbm, dsfont, xspace, cancel, mathtools). The delivered numbering is the unit's own. Assets: no retained span contains a figure environment, a graphics-inclusion command, a PDF-page inclusion, a table or a TikZ picture; no image file is copied into this package, the package declares no asset dependency and no image file is redistributed with it. Licence: version arXiv:2604.23918v1 of this article is distributed under the Creative Commons Attribution 4.0 International licence, as recorded in the arXiv licence metadata for this exact version and shown on the abstract page https://arxiv.org/abs/2604.23918v1. A full rights notice is delivered at licenses/arxiv-2604.23918-CC-BY-4.0.txt. Characters: every retained excerpt is pure ASCII, so the delivered engine, XeLaTeX with its default Latin Modern text font, reports no missing character.
\begin{align}
%%\label{merten}
%%\sum_{p\le x} \frac{1}{p} =& \log \log x + b_1+ O\Big(\frac{1}{\log x}\Big), \;  \\
%%\label{merten1}
%%\sum_{p\le x} \frac{\lambda^2(p)}{p} =& \log \log x + b_2+ O\Big(\frac{1}{\log x}\Big),  \\
\label{merten2-0}
\sum_{p\le x} \log p =& 
  x+O\Big(x \exp(-b_1\sqrt{\log x})\Big), \\
\label{merten2}
\sum_{p\le x} \chi_4(p)\log p =&  
  O\Big(x \exp(-b_2\sqrt{\log x})\Big),  \\
\label{sumpsigmalogp}
\sum_{p\le x} \frac{\log p}{p^{\sigma}} =& \int^{x}_1\frac {du}{u^{\sigma}}+ O\Big(1+x^{1-\sigma}\exp(-b_3\sqrt{\log x})\Big), \\
\label{sumpsigmalogpchi}
\sum_{p\le x} \frac{\chi_4(p)\log p}{p^{\sigma}} =& O\Big(1+x^{1-\sigma}\exp(-b_4\sqrt{\log x})\Big), \\
\label{merten4}
 \prod_{p\le x}(1-\frac {1}{p})^{-1}(1-\frac {\chi_4(p)}{p})^{-1}=& \frac {\pi}{4}e^{\gamma_0}\log x+O(1).
\end{align}

\begin{align}
\label{alphalower}
   \alpha_G(y^u,y)\geq &\frac 2{\log y}, \quad 1 \leq u \leq \frac {y}{8\log y}, \\
\label{alphausmall}
   \alpha_G(y^u,y) = & 1+O(\frac {1}{\log y}),  \quad 1 \leq u \leq 14, \\
\label{alphaupper}
   \alpha_G(y^u,y)\leq & 1-\frac {4}{\log y},  \quad u \geq 14, \\  
\label{equ: Saddle point approx in intro}
\alpha_G(y^u,y) = & 1 - \frac{\xi(u)}{\log y} + O \biggl( \frac{1}{(\log y)^2}+\frac {u}{y}\biggr), \quad 1 \leq u \leq \frac {y}{8\log y}, \\
\label{equ: Saddle point approx in intro1}
  \alpha_G(y^u,y)=&  1 - \frac{\log(u\log u)}{\log y} + O \biggl( \frac{\log\log u}{\log u\log y}+\frac{1}{(\log y)^2}+\frac {u}{y}\biggr), \quad 3 \leq u \leq \frac {y}{8\log y}.
\end{align}

\begin{align}
\label{alpharel}
\begin{split}
 \log x =\sum_{p \leq y} \Big(\frac {\log p}{p^{\alpha_G (x,y)}-1}+\frac {\chi_4(p)\log p}{p^{\alpha_G  (x,y)}-\chi_4(p)}\Big) =:\sum_{p \leq y} f_1(p^{\alpha_G  (x,y)}  (x,y), \chi_4(p))\log p, 
\end{split}
\end{align}     

\begin{align}
\label{sumpvy}
\begin{split}
  \sum_{v<p \leq y}\Big(\frac {\log p}{p^{\sigma}-1}+\frac {\chi_4(p)\log p}{p^{\sigma}-\chi_4(p)}\Big )=\sum_{v<p \leq y}\frac {(1+\chi_4(p)) \log p}{p^{\sigma}}+O(\sum_{v<p \leq y}\frac {\log p}{p^{2\sigma}}).
\end{split}
\end{align}   

\begin{align}
\label{vest}
\begin{split}
  v(\alpha_G)\log v(\alpha_G) \ll \frac {y^{1-\alpha_G}}{(1-\alpha_G)\log y}.
\end{split}
\end{align}   

\begin{align}
\label{sumpvy4}
\begin{split}
  u\log y =\int^{y}_1\frac {dw}{w^{\alpha_G}}+ O\Big(v(\alpha_G)\log v(\alpha_G)+y^{1-\alpha_G}\exp(-b_5\sqrt{\log y})\Big)+O(\sum_{p \leq y}\frac {\log p}{p^{2\alpha_G}}).
\end{split}
\end{align}  

\begin{lemma}
\label{lemphibound}
  With the notation as above. Let $y \geq 2, x=y^u$. We have uniformly for $1 \leq u \leq y/(8\log y), \alpha_G=\alpha_G(x,y) \geq 1/2+\varepsilon_1$ for any fixed $\varepsilon_1>0$ and $1 \leq k \leq 4$,
\begin{align}
\label{phikasymp}
  |\phi_k(\alpha_G,G;y)| \asymp u(\log y)^k. 
\end{align}
\end{lemma}  

\begin{proof}
  Note that we have for $1 \leq k \leq 4$,
\begin{align}
\label{phik}
\begin{split}
  |\phi_k(\alpha_G,G;y)| \asymp & \Big|\sum_{p \leq y}(\frac {p^{(k-1)\alpha_G}(\log p)^k}{(p^{\alpha_G}-1)^k}+\frac {\chi_4(p)p^{(k-1)\alpha_G}(\log p)^k}{(p^{\alpha_G}-\chi_4(p))^k}\Big |.
\end{split}
\end{align}     
   Now as $\alpha_G \geq 1/2+\varepsilon_1$, we have $2\alpha_G \geq 1+2\varepsilon_1$, so that by Taylor expansion, the right-hand side expression above is
\begin{align}
\label{phik1}
\begin{split}
  = & \Big|O(1)+\sum_{p \leq y}\frac {(1+\chi_4(p))(\log p)^k}{p^{\alpha_G}}\Big |. 
\end{split}
\end{align}     
  As $1+\chi_4(p) \geq 0$, we see that the above is
\begin{align}
\label{phik2}
\begin{split}
  \leq \Big |\sum_{p \leq y}\frac {(1+\chi_4(p))(\log p)^k}{p^{\alpha_G}}\Big |+O(1)
  \leq (\log y)^{k-1} \sum_{p \leq y}\frac {(1+\chi_4(p))\log p}{p^{\alpha_G}}+O(1) \ll u(\log y)^{k}, 
\end{split}
\end{align} 
 where the last estimation above follows from \eqref{alpharel} and \eqref{sumpvy} (by setting $v=1$ there), keeping in mind that we have $\alpha_G \geq 1/2+\varepsilon_1$ here.   
  
  We deduce from \eqref{phik}--\eqref{phik2} that
\begin{align}
\label{phikupper}
\begin{split}
  |\phi_k(\alpha_G,G;y)| \ll u(\log y)^{k}.
\end{split}
\end{align} 

  On the other hand, we infer from \eqref{sumpsigmalogp} that there exists a constant $b_6>0$ such that
\begin{align}
\label{phiklower1}
\begin{split}
  & \Big|O(1)+\sum_{p \leq y}\frac {(1+\chi_4(p))(\log p)^k}{p^{\alpha_G}}\Big | \\
\geq & \Big |\sum_{\sqrt{y} < p \leq y}\frac {(1+\chi_4(p))(\log p)^k}{p^{\alpha_G}}\Big |+O(1)
  \geq (\log y)^{k-1} \sum_{\sqrt{y} < p \leq y}\frac {(1+\chi_4(p))\log p}{p^{\alpha_G}}+O(1) \\
  \gg &  \frac {y^{1-\alpha_G}-y^{(1-\alpha_G)/2}}{1-\alpha_G}(\log y)^{k-1} +O((\log y)^{k-1} +(\log y)^{k-1}y^{1-\alpha_G}\exp(-b_6\sqrt{\log y})). 
\end{split}
\end{align}      
   Note that when $14 < u \leq y/(8\log y)$, we have by \eqref{alphaupper} that $\alpha_G \leq 1-4/\log y$. It is then easy to see that we have $\frac {y^{1-\alpha_G}-y^{(1-\alpha_G)/2}}{1-\alpha_G} \gg \frac {y^{1-\alpha_G}-1}{1-\alpha_G} \gg \frac {y^{1-\alpha_G}}{1-\alpha_G} $. On the other hand, we deduce from \eqref{vest}, \eqref{sumpvy4} and our assumption $\alpha_G \geq 1/2+\varepsilon_1$ that 
\begin{align*}
%%\label{phik4}
\begin{split}
 u \log y \ll \int^y_1\frac {dw}{w^{\alpha_G}}+O(1) \ll \frac {y^{1-\alpha_G}-1}{1-\alpha_G} \ll \frac {y^{1-\alpha_G}}{1-\alpha_G}. 
\end{split}
\end{align*}      

   We deduce from this and \eqref{phiklower1} that 
\begin{align*}
%%\label{phik4}
\begin{split}
  \Big|O(1)+\sum_{p \leq y}\frac {(1+\chi_4(p))(\log p)^k}{p^{\alpha_G}}\Big | 
  \gg &  \frac {y^{1-\alpha_G}}{1-\alpha_G}(\log y)^{k-1}+O((\log y)^{k-1}) \gg u(\log y)^{k}. 
\end{split}
\end{align*}      
    It follows from the above, \eqref{phik} and \eqref{phik1} that 
\begin{align}
\label{phiklower}
\begin{split}
  |\phi_k(\alpha_G,G;y)| \gg u(\log y)^{k}.
\end{split}
\end{align}    
  The above together with \eqref{phikupper} now implies \eqref{phikasymp} for the case $14 < u \leq y/(8\log y)$. 
  
  When $1 \leq u \leq 14$, we repeat the above arguments to see that it suffices to assume that $\alpha_G \geq 1-4/\log y$. Moreover, by \eqref{alphausmall}, we may also assume that $\alpha_G \leq 1+C/\log y$ for some constant $C>0$.  By \eqref{vest} and \eqref{sumpvy4}, we see that 
\begin{align}
\label{wintlower}
\begin{split}
  \int^y_1\frac {dw}{w^{\alpha_G}} \gg u \log y.
\end{split}
\end{align}   

  We now apply \eqref{sumpsigmalogp} and \eqref{sumpsigmalogpchi} to see that when $1-4/\log y \leq \alpha_G \leq 1+C/\log y$,  we have for any constant $0<c<1$, 
\begin{align}
\label{phik5}
\begin{split}
  \sum_{y^c < p \leq y}\frac {(1+\chi_4(p))\log p}{p^{\alpha_G}} =\int^y_{1}\frac {dw}{w^{\alpha_G}}+O(1+\int^{y^c}_1\frac {dw}{w^{\alpha_G}}).
\end{split}
\end{align}   
  
  As $\alpha_G \geq 1-4/\log y$, we have   
\begin{align*}
%%\label{phik7}
\begin{split}
 \int^{y^c}_1\frac {dw}{w^{\alpha_G}} \leq e^4\int^{y^c}_1\frac {dw}{w} \leq ce^4 \log y.
\end{split}
\end{align*}     
  Thus upon choosing $c$ small enough, we see from \eqref{wintlower}, \eqref{phik5} and the above that 
\begin{align*}
%%\label{phik7}
\begin{split}
  \sum_{y^c < p \leq y}\frac {(1+\chi_4(p))\log p}{p^{\alpha_G}} \gg u\log y. 
\end{split}
\end{align*}   
  We now repeat our arguments above for the case $u > 14$ to see that the estimation given in \eqref{phiklower} is valid for the case $1 \leq u \leq 14$ as well. Together with \eqref{phikupper}, it now implies \eqref{phikasymp} for the case  $1 \leq u \leq 14$. 
  This completes the proof of the lemma. 
\end{proof}

\end{document}
